\chapter{decision\_frame：把市场压成一行可回测状态}

\section{本章要解决的困惑}

为什么不直接在 fast backtest 里读 quote、trade、L2？因为 raw stream 是事件流，不是一行一行可直接判断 entry 的状态表。
\code{decision_frame_v1} 的作用，就是把市场事件压成策略在某个观察点能看到的一行。

\section{一行 decision frame 回答什么问题}

一行 \code{decision_frame} 可以粗略理解为：

\begin{quote}
在这个 L2/panel clock 时刻，策略能看到的 top-of-book、最近成交流、价格停滞状态和过去短路径状态是什么？
\end{quote}

关键字段如下：

\begin{longtable}{p{0.30\textwidth}p{0.24\textwidth}p{0.36\textwidth}}
\toprule
字段 & 类型 & 交易含义 \\
\midrule
\code{observed_seq} & int & 观察到的事件序号 \\
\code{local_ts_us} & int & 本地时间戳，微秒 \\
\code{best_bid_price} & float & 当前最好买价 \\
\code{best_ask_price} & float & 当前最好卖价 \\
\code{mid_price} & float & 中间价 \\
\code{spread_bps} & float & 顶层价差，bps \\
\code{trade_window_count} & int & 最近 trade window 内成交数量 \\
\code{trade_buy_amount} & float & 最近窗口 buy 主动量 \\
\code{trade_sell_amount} & float & 最近窗口 sell 主动量 \\
\code{trade_flow_imbalance} & float & TFI \\
\code{frames_since_mid_change} & float & mid 已经多少帧没变 \\
\code{past_event_25_bps} & float & 过去约 25 个事件 mid 变化 \\
\code{fwd_time_60s_bucket} & int & 用于限制同 bucket 重复触发 \\
\bottomrule
\end{longtable}

\warningline{\code{fwd_time_60s_bucket} 不是 future label；它是分桶 clock。不要被名字里的 fwd 吓到。}

\section{小例子：从盘口和成交到一行状态}

假设当前盘口：

\begin{lstlisting}
best_bid_price = 0.1120
best_bid_amount = 12000
best_ask_price = 0.1124
best_ask_amount = 9000
\end{lstlisting}

最近 2 秒成交：

\begin{lstlisting}
buy qty  = 500
sell qty = 0
\end{lstlisting}

过去 25 个事件 mid 几乎没动：

\begin{lstlisting}
past_event_25_bps = 0.2
frames_since_mid_change = 34
\end{lstlisting}

则：

\[
mid = 0.1122,
\quad
spread\_bps = \frac{0.1124 - 0.1120}{0.1122} \times 10000.
\]

\[
TFI = \frac{500 - 0}{500 + 0} = 1.
\]

这行状态就具备了旧策略判断 entry 所需的核心信息：

\begin{lstlisting}
TFI = 1
trade_window_count > 0
abs(past_event_25_bps) <= 0.5
frames_since_mid_change >= 25
spread_bps known
\end{lstlisting}

\section{decision frame 不是 label table}

这点必须讲清楚。一个合法的 runtime-safe decision frame 可以包含：

\begin{itemize}[leftmargin=2em]
  \item 当前 bid、ask、mid；
  \item 最近成交流；
  \item 当前或过去的 L2 状态；
  \item 当前时刻以前的 mid path 统计；
  \item 用于同频判断的 clock bucket。
\end{itemize}

它不能包含：

\begin{itemize}[leftmargin=2em]
  \item entry 后 60 秒收益；
  \item 未来最大涨幅 MFE；
  \item 未来最大回撤 MAE；
  \item 事后打分的 scored entries；
  \item 策略跑完后的 PnL。
\end{itemize}

如果把这些答案写进 decision frame，再让策略读它，就不是回测，是作弊。

\section{对应代码入口}

主要看：

\begin{lstlisting}
systems/ccusdt_replay_exchange/strategies/python/ccusdt_tfi_core_idle01/decision_frame.py
systems/ccusdt_replay_exchange/diagnostics/decision_frame_cache.py
\end{lstlisting}

\code{decision_frame.py} 里最值得先看：

\begin{lstlisting}
CACHE_FIELDS
DecisionFrameBuilder
DecisionL2Book
observe_trade
build_from_l2_updates
\end{lstlisting}

\code{decision_frame_cache.py} 里最值得先看：

\begin{lstlisting}
PARQUET_SCHEMA
cache_path
manifest_path
iter_cached_decision_frames
\end{lstlisting}

\section{手动执行检查点}

拿到一行 frame，先用纸算三件事：

\begin{enumerate}[leftmargin=2em]
  \item \code{spread_bps} 是否等于 \((ask-bid)/mid \times 10000\)？
  \item \code{trade_flow_imbalance} 的正负方向是否和 buy/sell 主动量一致？
  \item \code{past_event_25_bps} 小不小，\code{frames_since_mid_change} 高不高？
\end{enumerate}

能算出这三件事，就能进入 entry 章节。

